\editionPage[\pdfbookmark[0]{Corrigés des bilans}{bm-bilan-A-cor-auto-parent}]{bilan-A-cor-auto}{Bilan A · Corrigé}{Automatismes expliqués}{\hyperref[bilan-A-auto]{Énoncé p.\,\pageref*{bilan-A-auto}}}
\begin{blocEdition}{bilan-A-cor-auto}
\begin{questions}
\item Mêmes signes, produit positif : $\res{24}$.
\item Le produit est prioritaire : $7-10=\res{-3}$.
\item $6/8+1/8=\res{7/8}$.
\item L’inverse est $\res{-5/2}$ ; le produit des deux vaut 1.
\item $\res{4{,}6\times10^{-4}}$.
\item $11^2=121$ et $11\ge0$, donc $\res{11}$.
\item Distribuer 3 : $\res{6x-15}$.
\item $4x=16$, donc $\res{x=4}$.
\item $15\times80/100=\res{12}$.
\item $54/3{,}6=\res{15\ \mathrm{m/s}}$.
\item La troisième valeur est centrale : $\res8$.
\item $1-0{,}7=\res{0{,}3}$.\end{questions}
\end{blocEdition}
\editionPage{bilan-A-cor-1}{Bilan A · Corrigé 1}{Un atelier et ses achats}{\hyperref[bilan-A-ex-1]{Énoncé p.\,\pageref*{bilan-A-ex-1}}}
\begin{blocEdition}{bilan-A-cor-1}
\editionEtape{1 · Modéliser}Les prix sont $A=12+4n$ et $B=7n$ euros. Le premier tarif n’est pas proportionnel au nombre de séances à cause de la partie fixe.
\editionEtape{2 · Résoudre}$12+4n=7n$ donne $12=3n$, donc $\res{n=4}$. Les deux dépenses valent $\res{28}$ euros.
\editionEtape{3 · Appliquer la remise}Pour huit séances, A vaut $12+32=44$ euros. La remise est 4,40 euros, donc A coûte $\res{39{,}60}$ euros. B coûte $7\times8=56$ euros. A est moins cher de $\res{16{,}40}$ euros.
\end{blocEdition}
\editionPage{bilan-A-cor-2}{Bilan A · Corrigé 2}{Un terrain et un câble}{\hyperref[bilan-A-ex-2]{Énoncé p.\,\pageref*{bilan-A-ex-2}}}
\begin{blocEdition}{bilan-A-cor-2}
\editionEtape{1 · Pythagore}$BC^2=6^2+8^2=100$, donc $\res{BC=10\ \mathrm m}$. L’hypoténuse dépasse les deux côtés de l’angle droit.
\editionEtape{2 · Cosinus}$\cos\widehat{ABC}=AB/BC=6/10=0{,}6$. En degrés, $\widehat{ABC}=\arccos(0{,}6)\approx\res{53^\circ}$.
\editionEtape{3 · Les réserves}La longueur achetée est $10+2\times0{,}5=11$ m. Le coût vaut $11\times3{,}20=\res{35{,}20}$ euros. Les réserves s’ajoutent après le calcul de la distance.\label{fig-bilan-A-cor-2}\figureGeo[corrige]{bilan-a-geometrie}

\end{blocEdition}
\editionPage{bilan-A-cor-3}{Bilan A · Corrigé 3}{Une petite enquête}{\hyperref[bilan-A-ex-3]{Énoncé p.\,\pageref*{bilan-A-ex-3}}}
\begin{blocEdition}{bilan-A-cor-3}
\editionEtape{1 · Deux indicateurs}La somme est 205 min. La moyenne vaut $\res{205/12\ \mathrm{min}}\approx17{,}08$ min. Les rangs 6 et 7 sont 15 et 20 ; la médiane choisie vaut $\res{17{,}5}$ min.
\editionEtape{2 · Une part du total}Six élèves lisent au moins 20 min : fréquence $6/12=\res{1/2}$, soit 50\,\%. Le secteur correspondant mesure $\res{180^\circ}$.
\editionEtape{3 · Le hasard}Les six autres ont une durée inférieure à 20 min. Comme chaque élève a la même chance d’être choisi, la probabilité vaut $\res{1/2}$. On retrouve $1-1/2=1/2$ avec le contraire.
\end{blocEdition}
\editionPage{bilan-A-cor-4}{Bilan A · Corrigé 4}{Un cône et sa réduction}{\hyperref[bilan-A-ex-4]{Énoncé p.\,\pageref*{bilan-A-ex-4}}}
\begin{blocEdition}{bilan-A-cor-4}
\editionEtape{1 · Volume}$V=\pi\times3^2\times12/3=\res{36\pi\ \mathrm{cm}^3}\approx\res{113{,}1\ \mathrm{cm}^3}$.
\editionEtape{2 · Réduction}Le facteur de volume vaut $(1/2)^3=1/8$. Le petit cône a pour volume $36\pi/8=\res{4{,}5\pi\ \mathrm{cm}^3}$.
\editionEtape{3 · Contenance}Un cm$^3$ vaut un mL, donc la capacité est d’environ 113,1 mL. $\res{110\ \mathrm{mL}}$ peuvent y tenir dans ce modèle. Les dimensions intérieures sont supposées être celles données.\label{fig-bilan-A-cor-4}\figureGeo[corrige]{bilan-a-cone}

\end{blocEdition}
\editionPage{bilan-B-cor-auto}{Bilan B · Corrigé}{Automatismes expliqués}{\hyperref[bilan-B-auto]{Énoncé p.\,\pageref*{bilan-B-auto}}}
\begin{blocEdition}{bilan-B-cor-auto}
\begin{questions}
\item Signes opposés : $\res{-5}$.
\item Simplifier les facteurs : $\res{3/2}$.
\item $\res{84=2^2\times3\times7}$.
\item $\res{62\,000}$.
\item $81<90<100$, donc $\res{9<\sqrt{90}<10}$.
\item $\res{5x+5}$.
\item Facteur commun 6 : $\res{6(2x+3)}$.
\item $3x=12$, donc $\res{x=4}$.
\item $0{,}025\times1000=\res{25\ \mathrm L}$.
\item $12/8=\res{1{,}5}$.
\item $\res{3/10=0{,}3}$.
\item $(4-2;1+3)=\res{(2;4)}$.\end{questions}
\end{blocEdition}
\editionPage{bilan-B-cor-1}{Bilan B · Corrigé 1}{Un bassin et deux unités}{\hyperref[bilan-B-ex-1]{Énoncé p.\,\pageref*{bilan-B-ex-1}}}
\begin{blocEdition}{bilan-B-cor-1}
\editionEtape{1 · Convertir}$0{,}004$ m$^3$/s $=\res4$ L/s, soit $\res{240}$ L/min.
\editionEtape{2 · Durée}$t=V/q=720/240=\res3$ min. On suppose le débit constant jusqu’à la fin.
\editionEtape{3 · Volume restant}En 90 s, la pompe retire $4\times90=360$ L. Il reste $\res{360\ \mathrm L}$, soit $360/720=\res{1/2}$ de la capacité initiale. La durée choisie est bien antérieure à l’épuisement.
\end{blocEdition}
\editionPage{bilan-B-cor-2}{Bilan B · Corrigé 2}{Un parallélisme à prouver}{\hyperref[bilan-B-ex-2]{Énoncé p.\,\pageref*{bilan-B-ex-2}}}
\begin{blocEdition}{bilan-B-cor-2}
\editionEtape{1 · Réciproque}$AM/AB=4/10=2/5$ et $AN/AC=6/15=2/5$. Les points sont dans l’ordre sur les segments. La réciproque de Thalès donne $\res{(MN)\parallel(BC)}$.
\editionEtape{2 · Longueurs}Thalès donne $MN=(2/5)\times12=\res{4{,}8}$ cm. Puis $NC=AC-AN=15-6=\res9$ cm.
\editionEtape{3 · Aire}Les longueurs sont réduites au rapport $2/5$. Le rapport des aires est donc $(2/5)^2=\res{4/25}$, soit 16\,\%. Une aire n’est pas multipliée par le simple rapport des longueurs.\label{fig-bilan-B-cor-2}\figureGeo[corrige]{bilan-b-thales}

\end{blocEdition}
\editionPage{bilan-B-cor-3}{Bilan B · Corrigé 3}{Une série et un tirage}{\hyperref[bilan-B-ex-3]{Énoncé p.\,\pageref*{bilan-B-ex-3}}}
\begin{blocEdition}{bilan-B-cor-3}
\editionEtape{1 · Statistiques}La moyenne vaut $(6+16+24)/20=\res{2{,}3}$. Dans la série ordonnée, les rangs 10 et 11 portent tous deux 2 : médiane $\res2$.
\editionEtape{2 · Probabilités}Il y a $8+6=14$ cartes paires sur 20 : $P(\text{pair})=\res{7/10}$. Le contraire vaut $\res{3/10}$.
\editionEtape{3 · Modifier le total}Les 14 cartes paires ne changent pas. Pour représenter la moitié, elles doivent être dans un total de 28 cartes. On ajoute donc $28-20=\res8$ cartes portant 1. Vérification : $14/28=1/2$.
\end{blocEdition}
\editionPage{bilan-B-cor-4}{Bilan B · Corrigé 4}{Tripler une variable}{\hyperref[bilan-B-ex-4]{Énoncé p.\,\pageref*{bilan-B-ex-4}}}
\begin{blocEdition}{bilan-B-cor-4}
\editionEtape{1 · Suivi}À partir de 2, les tours donnent $\res{6,\ 18,\ 54,\ 162}$. Le premier 2 est une initialisation, pas un tour.
\editionEtape{2 · Puissance}Quatre multiplications par 3 donnent $2\times3^4=2\times81=\res{162}$.
\editionEtape{3 · Remonter le calcul}Avec une entrée $a$, le résultat vaut $81a$. L’équation $81a=405$ donne $\res{a=5}$. Contrôle : $5\times3^4=405$.\label{fig-bilan-B-cor-4}\figureGeo[corrige]{bilan-b-script}

\end{blocEdition}
